%! Sinusoidal Function Transformations — Unit 6, Lesson 6.7 — Warm-Up
\documentclass[10pt]{article}
\usepackage{precalculus-article}
\usepackage{precalculus-boxes}
\newcommand{\TallMath}[1]{$\displaystyle #1\rule[-1.4em]{0pt}{3.2em}$}

\begin{document}

\pageheader{Unit 6, Lesson 6.7}{Warm-Up}
\namedateperiod

\vspace{0.12in}

\begin{notesbox}{Spiral Review}

\textbf{1.}\quad Evaluate each expression using your knowledge of the unit circle.

\vspace{4pt}
\renewcommand{\arraystretch}{1.6}
\begin{tabular}{@{}p{0.22\linewidth}p{0.22\linewidth}p{0.22\linewidth}p{0.22\linewidth}@{}}
(a)\;$\sin\!\left(\dfrac{\pi}{2}\right) = $ \blank{1cm}
& (b)\;$\cos(\pi) = $ \blank{1cm}
& (c)\;$\sin\!\left(\dfrac{3\pi}{2}\right) = $ \blank{1cm}
& (d)\;$\cos(0) = $ \blank{1cm}
\end{tabular}

\vspace{0.14in}
\textbf{2.}\quad The graph of $y = \sin\theta$ is shown below with key features labeled.
Fill in each blank.

\vspace{4pt}
\begin{center}
\begin{tikzpicture}[x=1.1cm, y=1.1cm]
  \draw[linegray, very thin] (-0.2,-1.4) grid[xstep=0.5, ystep=0.5] (6.7,1.4);
  \draw[->] (-0.2,0) -- (6.8,0) node[right]{\small $\theta$};
  \draw[->] (0,-1.4) -- (0,1.5) node[above]{\small $y$};
  % Tick marks and labels on x-axis
  \foreach \x/\lbl in {1.571/$\frac{\pi}{2}$, 3.142/$\pi$, 4.712/$\frac{3\pi}{2}$, 6.283/$2\pi$} {
    \draw (\x,0.06) -- (\x,-0.06);
    \node[below, font=\scriptsize] at (\x,-0.06) {\lbl};
  }
  % Tick marks y-axis
  \foreach \y in {-1,1} {
    \draw (0.06,\y) -- (-0.06,\y);
    \node[left, font=\scriptsize] at (-0.06,\y) {$\y$};
  }
  % Sine curve
  \draw[navy, thick, domain=0:6.3, samples=100]
    plot (\x, {sin(deg(\x))});
\end{tikzpicture}
\end{center}

\vspace{4pt}
Amplitude $= $ \blank{1.5cm} \qquad Period $= $ \blank{1.5cm} \qquad Midline: $y = $ \blank{1.5cm}

\vspace{0.14in}
\textbf{3.}\quad A function $f$ is transformed by a vertical stretch of factor 4 and
a vertical shift up 2. If the parent function has the rule $f(\theta) = \sin\theta$,
write the rule of the transformed function.

\vspace{4pt}
Transformed function: $f(\theta) = $ \blank{4cm}

\vspace{0.14in}
\textbf{4.}\quad The function $g(\theta) = \sin(\theta - \pi)$ is a horizontal
translation of $y = \sin\theta$.

\vspace{4pt}
(a)\quad In which direction is the graph shifted? \writeline

\vspace{4pt}
(b)\quad By how many units? \writeline

\end{notesbox}

\end{document}
